\section{Desafios Epistemológicos e Estruturais}
\label{sec:desafios-educ}

A concretização da promessa tecnológica dos gêmeos digitais na academia odontológica esbarra em uma malha densa de desafios de ordem teórica, estrutural e comportamental que não podem ser subestimados ou simplificados sob a ótica do determinismo tecnológico.

O desafio cardeal incide sobre a \textbf{validação pedagógica e translação clínica}. Muito embora metanálises de metadados operacionais sugiram um incremento na curva de aprendizado, a literatura ainda anseia por Ensaios Clínicos Randomizados (ECRs) longitudinais, pragmáticos e multicêntricos que provem a transferência equívoca de habilidades adquiridas em ambiente sintético para a efetividade, resolutividade clínica e diminuição de taxas de iatrogenia a longo prazo no paciente humano~\cite{aisoc2026realising}. Há o risco epistemológico da simulação gerar "operadores de simulador" exímios, que contudo, careçam de competências não-técnicas como empatia, controle comportamental do paciente vivo e gestão da imprevisibilidade biológica.

\destaque{O hiato da translação ocorre quando o gêmeo digital é percebido pela academia meramente como um acessório de replicação mecânica, e não como um \textit{proxy} integrativo da fisiologia e do comportamento humano em rede.}

Substancialmente tangível é o \textbf{custo marginal de infraestrutura}. O dimensionamento financeiro de uma "Clínica Virtual" de alto padrão (exigindo hápticos com resistência 6-DOF, arranjos de servidores de GPU, licencas de software anuais) frequentemente suplanta o CAPEX (\textit{Capital Expenditure}) de laboratórios analógicos. Para ecossistemas acadêmicos em nações em desenvolvimento, essa barreira exige novos modelos de negócios, como a formação de consórcios interinstitucionais de compartilhamento de plataformas, e incentivo pesados no movimento \textit{open-source} e \textit{open-hardware} de simuladores, sob a pena de um apartheid tecnológico irreversível entre faculdades de elite e instituições públicas.

Em paralelo, a \textbf{inércia e resistência institucional} moldam um gargalo sociocultural agudo. Corpos docentes sêniores manifestam natural ceticismo quanto à abstração do ensino tradicionalmente empírico-artesanal. Superar a fricção não exige o banimento dos métodos clássicos, mas a modelagem de currículos mistos (\textit{blended}), onde o digital atua como filtro de pré-requisitos lógicos para o acesso humano. Esta mudança requer letramento digital massivo dos preceptores clínicos e incentivos reais de progressão acadêmica para a adoção de pedagogias baseadas em inteligência algorítmica.

Por fim, a resiliência da \textbf{infraestrutura de telecomunicações} impõe restrições severas. A promessa da arquitetura em nuvem descrita anteriormente é altamente volátil à estabilidade da largura de banda regional e ao fenômeno de \textit{jitter}. Na simulação cirúrgica, atrasos de sinalização da ordem de 100 milissegundos fragmentam a imersão sensorial e induzem síndromes de cinetose e dissociação visomotora, inviabilizando o aprendizado neural profundo das sinapses proprioceptivas.
